% https://www.crackap.com/ap/physics-c-electricity-and-magnetism/index.html
% https://www.apstudy.net/ap/physics-c-electricity-and-magnetism/frq-test1.html





Electric Potential Energy
(a) The electric potential energy for a point charge $\mathrm{q}_0$ in the electric field of a stationary point charge q , with a distance $r$ separating the charges is

$$
\mathrm{U}=\frac{1}{4 \pi \varepsilon_0} \frac{\mathrm{q}_0 \mathrm{q}}{\mathrm{r}}
$$


If electric potential energy at infinity is considered to be zero.
(b) Work done by electric force on a charge when it is moved from $A$ to $B$

$$
\mathrm{W}_{\mathrm{A} \rightarrow \mathrm{~B}}=\mathrm{U}_{\mathrm{A}}-\mathrm{U}_{\mathrm{B}}
$$


Electric Potential
(a) Potential is equal to potential energy per unit charge

$$
\mathrm{V}=\frac{\mathrm{U}}{\mathrm{q}_0}
$$


The potential for a point charge q at any point at a distance $r$ is

$$
\mathrm{V}=\frac{1}{4 \pi \varepsilon_0} \frac{\mathrm{q}}{\mathrm{r}}
$$


If potential at infinity is considered to be zero.









(b) Potential due to a collection of charges is the sum of the potentials due to each charge.

$$
\mathrm{V}_{\mathrm{n}}=\frac{1}{4 \pi \varepsilon_0} \sum_{\mathrm{i}}^{\mathrm{n}} \frac{\mathrm{q}_{\mathrm{i}}}{\mathrm{r}}
$$

(c) Potential due to a conducting sphere of radius r with charge $q$ (solid or hollow) at a distance $r$ from the centre

$$
\begin{array}{ll}
V=\left(\frac{1}{4 \pi \varepsilon_0}\right) \frac{q}{r} & \text { if }(r>R) \\
V=\left(\frac{1}{4 \pi \varepsilon_0}\right) \frac{q}{R} & \text { if }(r=R) \\
V=\frac{1}{4 \pi \varepsilon_0} \frac{q}{R} & \text { if }(r<R)
\end{array}
$$








(d) Relation between electric field and potential

$$
\begin{aligned}
& |\overrightarrow{\mathrm{E}}|=\left|-\frac{\partial \mathrm{V}}{\partial \ell}\right| \\
& =+\frac{\partial \mathrm{V} \mid}{\partial \ell}
\end{aligned}
$$


Electric Dipole Potential
(a) $\mathrm{V}=\frac{1}{4 \pi \varepsilon_0}\left(\frac{\mathrm{p} \cos \theta}{\mathrm{r}^2}\right)$
(b) Potential energy of a dipole in an external electric field $U(\theta)=-\vec{p} \cdot \vec{E}$





Capacitors
Capacitance of a parallel plate capacitor

$$
\mathrm{C}=\varepsilon_0 \frac{\mathrm{~A}}{\mathrm{~d}}
$$


Also, $C=\frac{Q}{V}$
Electric Field Energy
(a) $\mathrm{U}=\frac{1}{2} \mathrm{QV}=\frac{\mathrm{Q}^2}{2 \mathrm{C}}=\frac{1}{2} \mathrm{CV}^2$
(b) Energy density of energy stored in electric field $U=\frac{1}{2} \varepsilon_0 E^2$

Combination of Capacitor
(a) When capacitors are combined in series,

$$
\frac{1}{\mathrm{C}_{\mathrm{eq}}}=\frac{1}{\mathrm{C}_1}+\frac{1}{\mathrm{C}_2}+\frac{1}{\mathrm{C}_3}+
$$

$\qquad$
(b) When capacitors are connected in parallel,

$$
\mathrm{C}_{\mathrm{eq}}=\mathrm{C}_1+\mathrm{C}_2+\mathrm{C}_3+
$$

$\qquad$
(c) Spherical capacitor,

$$
\begin{aligned}
& \mathrm{C}=4 \pi \varepsilon_0 \frac{\mathrm{ab}}{\mathrm{~b}-\mathrm{a}}, \quad (When outer shell is earthed)\\
& \mathrm{C}=4 \pi \varepsilon_0 \frac{\mathrm{~b}^2}{\mathrm{~b}-\mathrm{a}}, \quad (When inner shell is earthed)\\
& $\mathrm{C}=4 \pi \varepsilon_0 \mathrm{R}$ (For a sphere of radius R )\\
\end{aligned}
$$





(d) Cylindrical capacitor, $\mathrm{C}=\frac{2 \pi \varepsilon_0 1}{\ln \left(\frac{\mathrm{~b}}{\mathrm{a}}\right)}$









Dielectrics
(a) Induced charge, $\mathrm{q}^{\prime}=\mathrm{q}\left(1-\frac{1}{\mathrm{~K}}\right)$
\begin{aligned}
&\text { (b) Polarisation } \mathrm{P}=\frac{\text { Dipole moment }}{\text { Volume }}\\
&\begin{aligned}
& \overrightarrow{\mathrm{P}}=\chi_0 \overrightarrow{\mathrm{E}} \\
& \frac{\chi_0}{\varepsilon_0}=\mathrm{K}-1
\end{aligned}
\end{aligned}